\documentclass[11pt,a4paper]{article}
\usepackage[utf8]{inputenc}
\usepackage[T1]{fontenc}
\usepackage[french]{babel}
\usepackage{amsmath, amssymb, amsthm}
\usepackage{geometry}
\geometry{margin=2.5cm}

\title{Corrigé - TD 1 : Dénombrement et Probabilités}
\author{MT 331 - INP Esisar / UGA}
\date{Année 2025-2026}

\begin{document}

\maketitle

[span_0](start_span)\section*{Exercice 1[span_0](end_span)}
Le tirage se fait simultanément, l'ordre n'a donc pas d'importance. Il s'agit de combinaisons. [span_1](start_span)On choisit 5 cartes parmi 32[span_1](end_span). Le nombre total d'issues possibles est $\binom{32}{5} = 201376$.

\subsection*{1. [span_2](start_span)Seulement des cartes noires[span_2](end_span)}
[span_3](start_span)\textbf{Raisonnement :} Un jeu de 32 cartes contient 16 cartes noires (8 trèfles et 8 piques)[span_3](end_span). Il faut choisir 5 cartes parmi ces 16.\\
[span_4](start_span)\textbf{Réponse :} $\binom{16}{5} = \frac{16!}{5!11!} = 4368$ façons[span_4](end_span).

\subsection*{2. [span_5](start_span)Seulement des figures ou as[span_5](end_span)}
[span_6](start_span)\textbf{Raisonnement :} Les "figures" précisées ici sont au nombre de 4 par couleur (valet, dame, roi, as), soit $4 \times 4 = 16$ cartes au total[span_6](end_span). On en choisit 5 parmi 16.\\
[span_7](start_span)\textbf{Réponse :} $\binom{16}{5} = 4368$ façons[span_7](end_span).

\subsection*{3. [span_8](start_span)Quatre as[span_8](end_span)}
[span_9](start_span)\textbf{Raisonnement :} Il faut choisir les 4 as parmi les 4 disponibles, puis choisir la dernière carte parmi les 28 cartes restantes du jeu[span_9](end_span).\\
[span_10](start_span)\textbf{Réponse :} $\binom{4}{4} \times \binom{28}{1} = 1 \times 28 = 28$ façons[span_10](end_span).

\subsection*{4. [span_11](start_span)Cinq figures, dont 3 noires[span_11](end_span)}
\textbf{Raisonnement :} Il y a 8 figures noires et 8 figures rouges. [span_12](start_span)Il faut choisir 3 figures parmi les 8 noires, et les 2 autres cartes doivent être des figures rouges (choisies parmi les 8) pour avoir exactement 5 figures dont 3 noires[span_12](end_span).\\
[span_13](start_span)\textbf{Réponse :} $\binom{8}{3} \times \binom{8}{2} = 56 \times 28 = 1568$ façons[span_13](end_span).

\subsection*{5. [span_14](start_span)Trois as, 2 dames et 1 seul carreau[span_14](end_span)}
\textbf{Raisonnement :} On tire 5 cartes au total (3 as + 2 dames). [span_15](start_span)L'unique carreau de la main doit obligatoirement être soit un as, soit une dame[span_15](end_span). On sépare en deux cas disjoints :
\begin{itemize}
    \item \textbf{Cas 1 : Le carreau est un as.} On choisit l'as de carreau (1 choix). Les 2 autres as sont choisis parmi les 3 as restants (non-carreau). Les 2 dames sont choisies parmi les 3 dames restantes (non-carreau). Soit $\binom{1}{1} \times \binom{3}{2} \times \binom{3}{2} = 1 \times 3 \times 3 = 9$.
    \item \textbf{Cas 2 : Le carreau est une dame.} On choisit la dame de carreau (1 choix). L'autre dame est choisie parmi les 3 restantes (non-carreau). Les 3 as sont choisis parmi les 3 as non-carreau. Soit $\binom{1}{1} \times \binom{3}{1} \times \binom{3}{3} = 1 \times 3 \times 1 = 3$.
\end{itemize}
[span_16](start_span)\textbf{Réponse :} $9 + 3 = 12$ façons[span_16](end_span).

\hrulefill

\section*{Exercice 2}

\subsection*{1. [span_17](start_span)Nombre total de comités[span_17](end_span)}
\textbf{Raisonnement :} Il y a au total $12 + 15 = 27$ personnes. [span_18](start_span)On choisit 6 membres sans distinction d'ordre[span_18](end_span).\\
[span_19](start_span)\textbf{Réponse :} $\binom{27}{6} = 296010$ comités possibles[span_19](end_span).

\subsection*{2. [span_20](start_span)Comités respectant la parité[span_20](end_span)}
[span_21](start_span)\textbf{Raisonnement :} La parité signifie 3 hommes et 3 femmes[span_21](end_span). On choisit 3 hommes parmi 12, et 3 femmes parmi 15.\\
[span_22](start_span)\textbf{Réponse :} $\binom{12}{3} \times \binom{15}{3} = 220 \times 455 = 100100$ comités[span_22](end_span).

\subsection*{3. [span_23](start_span)Condition de refus entre A et B[span_23](end_span)}
[span_24](start_span)\textbf{Raisonnement :} Madame A et Monsieur B ne doivent pas siéger ensemble[span_24](end_span). On compte le nombre de comités paritaires où ils siègent \textit{tous les deux}, et on le soustrait au total paritaire calculé précédemment.
[span_25](start_span)S'ils y sont tous les deux, il reste à choisir 2 femmes parmi 14 (pour compléter avec Mme A) et 2 hommes parmi 11 (pour compléter avec M B)[span_25](end_span).\\
[span_26](start_span)\textbf{Réponse :} $100100 - \left( \binom{14}{2} \times \binom{11}{2} \right) = 100100 - (91 \times 55) = 100100 - 5005 = 95095$ comités[span_26](end_span).

\subsection*{4. [span_27](start_span)Désignation du bureau[span_27](end_span)}
[span_28](start_span)\textbf{Raisonnement :} Parmi les 6 membres, on désigne un président, un secrétaire et un trésorier (3 postes distincts)[span_28](end_span). L'ordre compte, c'est un arrangement de 3 personnes parmi 6.\\
[span_29](start_span)\textbf{Réponse :} $A_6^3 = \frac{6!}{(6-3)!} = 6 \times 5 \times 4 = 120$ bureaux possibles[span_29](end_span).

\hrulefill

[span_30](start_span)\section*{Exercice 3[span_30](end_span)}
[span_31](start_span)\textbf{Raisonnement :} On utilise l'événement contraire : "les $n$ personnes ont toutes des jours d'anniversaire différents"[span_31](end_span). [span_32](start_span)Il y a $365^n$ cas possibles au total (nombre d'applications de l'ensemble des personnes vers les jours de l'année, sans années bissextiles[span_32](end_span)). Le nombre de cas favorables à l'événement contraire est le nombre d'arrangements de $n$ jours parmi 365, soit $A_{365}^n$.\\
[span_33](start_span)\textbf{Réponse :} La probabilité est $P = 1 - \frac{A_{365}^n}{365^n}$[span_33](end_span).

\hrulefill

[span_34](start_span)\section*{Exercice 4[span_34](end_span)}

\subsection*{1. [span_35](start_span)Maths groupés et Physique groupés[span_35](end_span)}
\textbf{Raisonnement :} Il y a 2 "blocs" à ordonner : le bloc Maths et le bloc Physique (2! façons de les disposer). [span_36](start_span)À l'intérieur du bloc Maths, il y a 12! façons de ranger les livres[span_36](end_span). À l'intérieur du bloc Physique, il y a 10! façons.\\
\textbf{Réponse :} $2! \times 12! [span_37](start_span)\times 10!$ rangements possibles[span_37](end_span).

\subsection*{2. [span_38](start_span)Seulement les Maths groupés[span_38](end_span)}
\textbf{Raisonnement :} On considère les 12 livres de maths comme un seul gros bloc. [span_39](start_span)Ce bloc s'ajoute aux 10 livres de physique, soit 11 éléments à permuter ($11!$ façons)[span_39](end_span). À l'intérieur du bloc des maths, les 12 livres peuvent être permutés de $12!$ façons.\\
\textbf{Réponse :} $11! [span_40](start_span)\times 12!$ rangements possibles[span_40](end_span).

\hrulefill

[span_41](start_span)\section*{Exercice 5[span_41](end_span)}
[span_42](start_span)\textbf{Raisonnement global :} L'armoire contient 20 gants au total (10 paires distinctes)[span_42](end_span). [span_43](start_span)Le nombre de tirages possibles de 4 gants est $\binom{20}{4} = 4845$[span_43](end_span).

[span_44](start_span)\subsection*{(i) Deux paires assorties[span_44](end_span)}
\textbf{Raisonnement :} Il faut tirer 2 paires complètes. [span_45](start_span)Il suffit de choisir 2 paires parmi les 10 paires disponibles[span_45](end_span).\\
\textbf{Réponse :} Nombre de cas favorables : $\binom{10}{2} = 45$. Probabilité = $\frac{45}{4845} = \frac{3}{323}$.

[span_46](start_span)\subsection*{(iii) Une paire et une seule[span_46](end_span)}
\textbf{Raisonnement :} On choisit la paire assortie parmi 10 ($\binom{10}{1}$). Pour les 2 autres gants, il faut qu'ils proviennent de 2 paires différentes. [span_47](start_span)On choisit 2 paires parmi les 9 restantes ($\binom{9}{2}$) et pour chacune de ces deux paires, on choisit le gant gauche ou droit ($2^2$ choix)[span_47](end_span).\\
\textbf{Réponse :} Nombre de cas = $10 \times 36 \times 4 = 1440$. Probabilité = $\frac{1440}{4845} = \frac{96}{323}$.

[span_48](start_span)\subsection*{(ii) Au moins une paire[span_48](end_span)}
\textbf{Raisonnement :} "Au moins une paire" signifie soit "une paire exacte", soit "deux paires". [span_49](start_span)Les deux événements sont incompatibles, on additionne leurs probabilités[span_49](end_span).\\
\textbf{Réponse :} Probabilité = $\frac{45}{4845} + \frac{1440}{4845} = \frac{1485}{4845} = \frac{99}{323}$.

\hrulefill

[span_50](start_span)\section*{Exercice 6[span_50](end_span)}
[span_51](start_span)\textbf{Raisonnement :} C'est un tirage simultané (ou sans remise) de $r$ boules parmi $n$[span_51](end_span). Le nombre total de tirages est $\binom{n}{r}$. [span_52](start_span)[span_53](start_span)Pour obtenir exactement $k$ boules noires, on doit choisir $k$ boules parmi les $m$ noires, et les $r-k$ boules restantes parmi les $n-m$ boules blanches[span_52](end_span)[span_53](end_span).\\
[span_54](start_span)\textbf{Réponse :} La probabilité est $P = \frac{\binom{m}{k} \times \binom{n-m}{r-k}}{\binom{n}{r}}$ (Loi hypergéométrique)[span_54](end_span).

\hrulefill

[span_55](start_span)\section*{Exercice 7[span_55](end_span)}
[span_56](start_span)\textbf{Raisonnement global :} Nombre total de mains de 5 cartes parmi 32 : $\binom{32}{5} = 201376$[span_56](end_span). Un jeu de 32 a 8 hauteurs distinctes.

\begin{itemize}
    \item \textbf{(i) Une seule paire :} On choisit la hauteur de la paire ($\binom{8}{1}$), puis 2 cartes de cette hauteur ($\binom{4}{2}$). On choisit ensuite 3 hauteurs distinctes parmi les 7 restantes ($\binom{7}{3}$), puis 1 carte pour chacune d'elles ($4^3$).\\
    \textbf{Probabilité :} $\frac{\binom{8}{1}\binom{4}{2}\binom{7}{3} \times 4^3}{201376} = \frac{107520}{201376} \approx 0.5339$.
    
    \item \textbf{(ii) Deux paires :} On choisit 2 hauteurs pour les paires ($\binom{8}{2}$), 2 cartes par hauteur ($\binom{4}{2}^2$). La cinquième carte est d'une autre hauteur ($\binom{6}{1} \times \binom{4}{1}$).\\
    \textbf{Probabilité :} $\frac{\binom{8}{2}\binom{4}{2}^2 \times 6 \times 4}{201376} = \frac{24192}{201376} \approx 0.1201$.
    
    \item \textbf{(iii) Un brelan :} On choisit la hauteur du brelan ($\binom{8}{1}$) et 3 cartes ($\binom{4}{3}$). Pour éviter un full, on prend 2 hauteurs distinctes pour les 2 cartes restantes ($\binom{7}{2}$) et 1 carte de chaque ($4^2$).\\
    \textbf{Probabilité :} $\frac{\binom{8}{1}\binom{4}{3}\binom{7}{2} \times 4^2}{201376} = \frac{10752}{201376} \approx 0.0534$.
    
    \item \textbf{(iv) Un carré :} On choisit la hauteur du carré ($\binom{8}{1}$), on prend les 4 cartes ($\binom{4}{4} = 1$), et 1 carte parmi les 28 autres.\\
    \textbf{Probabilité :} $\frac{\binom{8}{1} \times 1 \times 28}{201376} = \frac{224}{201376} \approx 0.0011$.
\end{itemize}

\hrulefill

[span_57](start_span)\section*{Exercice 8[span_57](end_span)}
[span_58](start_span)Composition de ANNIVERSAIRE (12 lettres) : A(2), N(2), I(2), V(1), E(2), R(2), S(1)[span_58](end_span).

\subsection*{1. [span_59](start_span)Nombre d'anagrammes[span_59](end_span)}
\textbf{Raisonnement :} Il y a 12 lettres avec des répétitions (5 paires de lettres identiques). [span_60](start_span)C'est une permutation avec répétition[span_60](end_span).\\
[span_61](start_span)\textbf{Réponse :} $\frac{12!}{2!2!2!2!2!} = \frac{12!}{32} = 14968800$ mots[span_61](end_span).

\subsection*{2. [span_62](start_span)Commençant et finissant par une voyelle[span_62](end_span)}
\textbf{Raisonnement :} Les voyelles sont A, I, E (chacune 2 fois, soit 6 voyelles). Il y a deux cas pour le choix des lettres aux extrémités :
\begin{itemize}
    \item Même voyelle aux deux bouts (ex: A...A) : 3 choix pour la voyelle. Les 10 lettres restantes ont 4 paires. Permutations : $3 \times \frac{10!}{2!2!2!2!} = 3 \times 226800 = 680400$.
    \item Voyelles différentes aux bouts (ex: A...E) : $A_3^2 = 6$ choix pour le couple ordonné. Les 10 lettres restantes conservent 5 paires. Permutations : $6 \times \frac{10!}{2!2!2!2!2!} = 6 \times 113400 = 680400$.
\end{itemize}
[span_63](start_span)\textbf{Réponse :} $680400 + 680400 = 1360800$ mots[span_63](end_span).

\subsection*{3. [span_64](start_span)Toutes les voyelles groupées[span_64](end_span)}
\textbf{Raisonnement :} On lie les 6 voyelles en 1 seul bloc. Ce bloc + les 6 consonnes forment 7 éléments à permuter, avec les paires de N et R ($\frac{7!}{2!2!}$). [span_65](start_span)À l'intérieur du bloc voyelles, on permute les 6 lettres avec les paires de A, I, E ($\frac{6!}{2!2!2!}$)[span_65](end_span).\\
[span_66](start_span)\textbf{Réponse :} $\frac{7!}{4} \times \frac{6!}{8} = 1260 \times 90 = 113400$ mots[span_66](end_span).

\hrulefill

[span_67](start_span)\section*{Exercice 9[span_67](end_span)}

\subsection*{1. [span_68](start_span)Tirage simultané de 5 boules[span_68](end_span)}
\textbf{(a)} $\Omega$ est l'ensemble des combinaisons de 5 boules parmi 16 (les boules étant discernables). [span_69](start_span)$\text{Card}(\Omega) = \binom{16}{5} = 4368$[span_69](end_span).\\
[span_70](start_span)\textbf{(b)} $P(2B, 3N) = \frac{\binom{6}{2} \times \binom{10}{3}}{\binom{16}{5}}$[span_70](end_span).

\subsection*{2. [span_71](start_span)Tirage successif de 2 boules, sans remise[span_71](end_span)}
[span_72](start_span)\textbf{(a) Classé :} $\Omega$ est l'ensemble des arrangements de 2 boules parmi 16. $\text{Card}(\Omega) = A_{16}^2 = 16 \times 15 = 240$[span_72](end_span).\\
[span_73](start_span)\textbf{(b)} $P(\text{1B puis 1N}) = \frac{6 \times 10}{240} = \frac{1}{4}$[span_73](end_span).\\
\textbf{(c) Non classé :} $\Omega$ est l'ensemble des combinaisons. $\text{Card}(\Omega) = \binom{16}{2} = 120$. [span_74](start_span)$P(\text{1B et 1N}) = \frac{\binom{6}{1} \binom{10}{1}}{120} = \frac{60}{120} = \frac{1}{2}$[span_74](end_span).

\subsection*{3. [span_75](start_span)Tirage avec remise de 2 boules[span_75](end_span)}
[span_76](start_span)\textbf{(a) Classé :} $\Omega$ est l'ensemble des listes de 2 boules parmi 16. $\text{Card}(\Omega) = 16^2 = 256$[span_76](end_span).\\
[span_77](start_span)\textbf{(b)} $P(\text{1B puis 1N}) = \frac{6 \times 10}{256} = \frac{60}{256} = \frac{15}{64}$[span_77](end_span).\\
\textbf{(c) Non classé :} Il n'y a plus d'ordre imposé (la boule noire peut être la première ou la deuxième). [span_78](start_span)On a donc $P = 2 \times P(\text{1B puis 1N}) = 2 \times \frac{15}{64} = \frac{15}{32}$[span_78](end_span).

\hrulefill

[span_79](start_span)\section*{Exercice 10[span_79](end_span)}
[span_80](start_span)\textbf{Raisonnement :} On calcule la probabilité de l'événement contraire pour chaque scénario[span_80](end_span).
\begin{itemize}
    [span_81](start_span)\item Au moins un 6 en 4 lancers : $P_1 = 1 - \left(\frac{5}{6}\right)^4 \approx 1 - 0.482 = 0.518$[span_81](end_span).
    \item Au moins un double 6 en 24 lancers de 2 dés : La probabilité de ne pas faire un double 6 est de 35/36. [span_82](start_span)Donc $P_2 = 1 - \left(\frac{35}{36}\right)^{24} \approx 1 - 0.509 = 0.491$[span_82](end_span).
\end{itemize}
[span_83](start_span)\textbf{Réponse :} Il est plus probable d'obtenir au moins un 6 en lançant 4 fois de suite un dé[span_83](end_span).

\hrulefill

[span_84](start_span)\section*{Exercice 11[span_84](end_span)}
\textbf{Raisonnement :} Pour gagner, le joueur 2 doit remporter les 2 parties suivantes (probabilité de $\frac{1}{2} \times \frac{1}{2} = \frac{1}{4}$). [span_85](start_span)Dans tous les autres cas (le joueur 1 gagne la partie 4, ou le joueur 2 gagne la 4 et le joueur 1 gagne la 5), le joueur 1 emporte la mise (probabilité de $1 - \frac{1}{4} = \frac{3}{4}$)[span_85](end_span).\\
[span_86](start_span)\textbf{Réponse :} La mise doit être répartie proportionnellement à leurs chances de victoire : $\frac{3}{4}$ de la mise pour le premier joueur et $\frac{1}{4}$ pour le second[span_86](end_span).

\hrulefill

[span_87](start_span)\section*{Exercice 12[span_87](end_span)}
\subsection*{1. [span_88](start_span)Choix possibles[span_88](end_span)}
[span_89](start_span)\textbf{Réponse :} Il faut choisir 4 voitures parmi 20, soit $\binom{20}{4} = 4845$ choix[span_89](end_span).

\subsection*{2. [span_90](start_span)Au moins une voiture en panne[span_90](end_span)}
[span_91](start_span)\textbf{Raisonnement :} Il y a 5 voitures en panne et 15 qui fonctionnent[span_91](end_span).
\begin{itemize}
    \item \textbf{Façon 1 (Événement contraire) :} $1 - P(\text{0 panne}) = 1 - \frac{\binom{15}{4}}{\binom{20}{4}} = 1 - \frac{1365}{4845} \approx 0.718$.
    [span_92](start_span)\item \textbf{Façon 2 (Somme directe) :} $\frac{\binom{5}{1}\binom{15}{3} + \binom{5}{2}\binom{15}{2} + \binom{5}{3}\binom{15}{1} + \binom{5}{4}\binom{15}{0}}{\binom{20}{4}}$[span_92](end_span).
\end{itemize}

\subsection*{3. [span_93](start_span)Exactement 2 voitures en panne[span_93](end_span)}
[span_94](start_span)\textbf{Réponse :} On choisit 2 en panne parmi 5, et 2 qui fonctionnent parmi 15 : $\frac{\binom{5}{2} \times \binom{15}{2}}{\binom{20}{4}} = \frac{10 \times 105}{4845} = \frac{1050}{4845} \approx 0.217$[span_94](end_span).

\hrulefill

[span_95](start_span)\section*{Exercice 13[span_95](end_span)}
\textbf{Raisonnement :} C'est le nombre de partitions ordonnées d'un ensemble de 4 personnes (les nombres de Fubini). [span_96](start_span)On sépare selon le nombre $k$ de places distinctes à l'arrivée (les ex-aequo forment des "blocs")[span_96](end_span):
\begin{itemize}
    \item 1 bloc (tous 1er ex-aequo) : 1 façon.
    \item 2 blocs : Les partitions possibles sont du type "3 et 1" (4 façons) ou "2 et 2" (3 façons), total 7 partitions. Pour chacune, 2! ordres possibles d'arrivée. $7 \times 2 = 14$ façons.
    \item 3 blocs : Partitions "2, 1 et 1" (6 partitions). Ordonnées : $6 \times 3! = 36$ façons.
    \item 4 blocs (aucun ex-aequo) : $4! = 24$ façons.
\end{itemize}
[span_97](start_span)\textbf{Réponse :} $1 + 14 + 36 + 24 = 75$ classements possibles[span_97](end_span).

\hrulefill

[span_98](start_span)\section*{Exercice 14[span_98](end_span)}
[span_99](start_span)\textbf{Raisonnement :} Bien qu'il y ait autant de combinaisons brutes donnant une somme de 9 (126, 135, 144, 225, 234, 333) que donnant 10 (136, 145, 226, 235, 244, 334), il faut tenir compte des permutations, les dés étant physiquement distincts[span_99](end_span). 
\begin{itemize}
    \item Pour la somme de 9, les permutations donnent : $6 + 6 + 3 + 3 + 6 + 1 = 25$ façons.
    \item Pour la somme de 10, les permutations donnent : $6 + 6 + 3 + 6 + 3 + 3 = 27$ façons.
\end{itemize}
[span_100](start_span)\textbf{Réponse :} Il y a 27 cas favorables pour obtenir 10 contre 25 cas pour obtenir 9. Le total de 10 est donc plus probable[span_100](end_span).

\end{document}
